% Resume - Vivekananda Swamy Mattam
% Optimized for Autonomous Mobile Robots roles

\documentclass[10pt,letterpaper]{article}
\usepackage[utf8]{inputenc}
\usepackage[margin=0.4in, top=0.25in, bottom=0.25in]{geometry}
\usepackage{enumitem}
\usepackage{hyperref}
\usepackage{titlesec}
\usepackage{xcolor}

% Link colors - blue
\hypersetup{
    colorlinks=true,
    linkcolor=blue,
    urlcolor=blue
}

% Formatting
\pagestyle{empty}
\setlength{\parindent}{0pt}
\setlength{\parskip}{0pt}

% Section formatting - bold uppercase with line
\titleformat{\section}{\normalsize\bfseries\uppercase}{}{0em}{}[\vspace{-4pt}\rule{\textwidth}{0.4pt}]
\titlespacing{\section}{0pt}{5pt}{2pt}

% Tight lists
\setlist[itemize]{leftmargin=0.15in, itemsep=0pt, parsep=0pt, topsep=1pt}

\begin{document}

%----------HEADING----------
\begin{center}
    {\Large \textbf{VIVEKANANDA SWAMY MATTAM}} \\
    \vspace{4pt}
    \small
    +1 (949) 400-7126 \quad $\bullet$ \quad New York, USA \\
    \href{mailto:vm2677@nyu.edu}{vm2677@nyu.edu} \quad $\bullet$ \quad
    \href{https://linkedin.com/in/vivek-mattam-a8590b23a}{linkedin} \quad $\bullet$ \quad
    \href{https://vivekmattam02.github.io}{vivekmattam02.github.io}
\end{center}
\vspace{-4pt}

%----------EDUCATION----------
\section{Education}

\textbf{Masters of Science - Mechatronics and Robotics Engineering} \hfill Sep 2024 - May 2026 \\
\textit{New York University - Tandon School of Engineering} \\[2pt]
\textbf{B.Tech. Mechanical Engineering} \hfill Sep 2019 - May 2023 \\
\textit{VNR VJIET, Hyderabad, India}
\vspace{-2pt}

%----------EXPERIENCE----------
\section{Experience}

\href{https://wp.nyu.edu/aarab/}{\textbf{Course Assistant - Autonomous Mobile Robots}} $|$ \textit{NYU Tandon, Prof. Aliasghar Arab} \hfill Sep 2024 - Present
\begin{itemize}
    \item Sole CA for graduate robotics course; developed 9 lectures, 5 assignments, and \href{https://vivekmattam02.github.io/amr-notes/}{course website} covering MPC, Control Barrier Functions, and motion planning.
\end{itemize}
\vspace{1pt}

\textbf{Robotics Engineering Intern} \hfill Feb 2024 - Aug 2024 \\
\textit{Xmachines - Agricultural Robotics Startup}
\begin{itemize}
    \item Built Flask server for real-time sensor streaming integrating MPU-6050 IMU and Arducam IMX219 camera feeds for robot state monitoring during field tests.
    \item Developed prototype object detection tunnel for weeder robot; automated size classification and routing for sorting process.
\end{itemize}

%----------PROJECTS----------
\section{Projects}

\href{https://github.com/vivekmattam02/ros2-ackermann-racing-navigation}{\textcolor{blue}{\textbf{Vortex: Autonomous RC Car with Edge-Deployed Distilled Perception}}} $|$ \textit{Bell Labs Funded, MS Project} \hfill Jan 2025 - Present
\begin{itemize}
    \item Compressed a 900M-parameter foundation model ensemble (SAM2-Large + Depth Anything V2-Large) into a 3M-parameter MobileNetV3 student via multi-task knowledge distillation, achieving 300x parameter reduction while retaining real-time depth and segmentation.
    \item Deployed multi-task perception on Jetson Orin Nano at \textasciitilde5ms/frame (200 FPS) using TensorRT FP16 within a 15W power envelope, detecting obstacles a 2D LiDAR cannot see (glass doors, overhanging objects, chairs).
    \item Designed multi-sensor fusion pipeline integrating RPLiDAR A2M8, Orbbec Femto Bolt (ToF depth + 200Hz IMU + RGB), and dual magnetic encoders with EKF localization, fusing learned monocular depth with hardware ToF to fill coverage gaps.
    \item Architected full autonomy stack in ROS 2 Humble: 6 custom nodes (VESC motor control, encoder odometry, TensorRT inference, depth fusion, class-aware costmap, stamper bridge), Nav2 with MPPI controller, SLAM Toolbox, and Gazebo simulation.
\end{itemize}
\vspace{1pt}

\href{https://github.com/vivekmattam02/spot}{\textcolor{blue}{\textbf{Autonomous Person Following on Boston Dynamics Spot}}} \hfill Jan 2025 - Present
\begin{itemize}
    \item Built real-time person-following system using visual servoing with YOLOv8 detection, computing lateral, distance, and pitch errors for smooth 10Hz velocity commands.
    \item Implemented body pitch control for tracking people on stairs and slopes; exploring ViNT and GNM for goal-conditioned predictive navigation.
\end{itemize}
\vspace{1pt}

\href{https://github.com/vivekmattam02/LAC_26}{\textcolor{blue}{\textbf{Lunar Autonomy Challenge - Perception-Aware Navigation}}} \hfill May 2024 - Nov 2024
\begin{itemize}
    \item Built autonomous navigation for simulated lunar rover without GPS or prior maps, combining ORB-SLAM3 for visual odometry with FoundationStereo for depth estimation.
    \item Developed perception-aware costmap that propagates depth uncertainty through the pipeline; planner actively avoids shadowed craters and texture-poor surfaces where perception is unreliable.
    \item Implemented A* planning with uncertainty-weighted costs and pure pursuit control for the full autonomy loop.
\end{itemize}
\vspace{1pt}

\textbf{Visual Navigation \& Perception} $|$ \textit{\href{https://ai4ce.github.io/CityWalker/}{AI4CE Lab} / VIP Self-Drive, NYU} \hfill 2024 - Present
\begin{itemize}
    \item Deployed \href{https://ai4ce.github.io/CityWalker/}{CityWalker} (CVPR 2025) on FrodoBots EarthRover under \href{https://scholar.google.com/citations?user=YeG8ZM0AAAAJ}{Prof. Chen Feng}; developed Depth Barrier Regularization (DBR) for safety-aware training using monocular depth, 15\% reduction in depth violations (targeting CoRL 2026).
    \item Built camera-only TurtleBot3 navigation using ORB visual SLAM with A* planning; maze solving with CosPlace + SuperGlue achieving 79\% inlier ratio on pose verification.
    \item Perception pipeline from scratch: RANSAC plane fitting (66\% inliers on 113K points), ICP alignment, YOLOv11+ByteTrack tracking across 12,700+ images.
\end{itemize}
\vspace{1pt}

\href{https://github.com/Nishant-ZFYII/rob6323_go2_project}{\textcolor{blue}{\textbf{Reinforcement Learning for Quadruped Locomotion}}} $|$ \textit{NVIDIA Isaac Lab} \hfill Fall 2024
\begin{itemize}
    \item Trained Unitree Go2 with PPO implementing reward shaping, Raibert heuristic gait coordination, and actuator friction modeling with domain randomization.
    \item Achieved 2x baseline velocity tracking on flat and rough terrain for sim-to-real transfer.
\end{itemize}

%----------SKILLS----------
\section{Skills}

\begin{itemize}[leftmargin=0in, label={}, itemsep=0pt]
    \item \textbf{Languages:} Python, C++, C, MATLAB
    \item \textbf{Autonomous Navigation:} ROS/ROS 2, Nav2, SLAM Toolbox, ORB-SLAM3, Visual Odometry, EKF, MPC, MPPI, A*, Pure Pursuit
    \item \textbf{Perception \& ML:} PyTorch, OpenCV, CosPlace, SuperGlue, YOLOv8/v11, ByteTrack, RANSAC, ICP, Open3D
    \item \textbf{Hardware \& Simulation:} Gazebo, Isaac Lab, Jetson Orin Nano, Arduino, TurtleBot3, Boston Dynamics Spot, Docker, Git, Linux
\end{itemize}

\end{document}
